% Challenge dossier: VI.01.C3
\subsection*{Challenge VI.01.C3 --- Why the coefficient field matters to the union formula}
\label{chal:vi01-c3}

\paragraph{Challenge.}
Suppose one tries to define zero loci in \(R^n\) for a commutative ring \(R\) with
identity, using the same evaluation definition.  Prove that
\[
V(IJ)=V(I)\cup V(J)
\]
for all polynomial ideals \(I,J\subseteq R[x_1,\dots,x_n]\) if \(R\) is an integral
domain.  Show that the statement can fail when \(R\) has zero divisors by taking
\[
R=\mathbb Z/6\mathbb Z,
\qquad I=(2),
\qquad J=(3).
\]
In fact, show that validity of the formula for all constant principal ideals forces \(R\)
to be a domain.

\ifdefined\FullProblemDossiers
\paragraph{Why this challenge matters.}
The familiar finite-union proof silently uses the absence of zero divisors in the value
ring.  This challenge identifies the exact algebraic hypothesis behind that step.

\paragraph{Guided hints.}
In a domain, if \(f(a)g(a)=0\), then one factor is zero.  For \(\mathbb Z/6\mathbb Z\), note
that \(2\cdot3=0\) although neither constant is zero.

\paragraph{Solution.}
Assume \(R\) is a domain.  The inclusion
\[
V(I)\cup V(J)\subseteq V(IJ)
\]
is immediate.  Conversely, if \(a\notin V(I)\cup V(J)\), choose \(f\in I\), \(g\in J\)
with \(f(a)\neq0\), \(g(a)\neq0\).  Since \(R\) is a domain,
\[
f(a)g(a)\neq0,
\]
so \(a\notin V(IJ)\).  Hence equality holds.

Now take \(R=\mathbb Z/6\mathbb Z\).  Since the constant polynomial \(2\) never evaluates
to \(0\),
\[
V((2))=\varnothing,
\]
and likewise \(V((3))=\varnothing\).  But
\[
(2)(3)=(6)=(0),
\]
so
\[
V((2)(3))=V((0))=R^n,
\]
which is not the union of the two empty sets.

Finally, suppose the formula holds for all constant principal ideals.  If nonzero
\(a,b\in R\) satisfied \(ab=0\), then
\[
V((a))=V((b))=\varnothing
\]
but
\[
V((ab))=V((0))=R^n,
\]
a contradiction.  Thus \(R\) has no nonzero zero divisors and is a domain.

\paragraph{Extensions.}
Explore what the correct geometric replacement is over arbitrary rings; this leads toward
prime spectra and schemes later in Volume VI.

\paragraph{Provenance.}
New challenge explicitly exposing the field/domain hypothesis hidden in the classical
finite-union identity.
\fi
